\documentclass[11pt,a4paper]{article}

%====================================================
% PACKAGES
%====================================================
\usepackage[utf8]{inputenc}
\usepackage[T1]{fontenc}
\usepackage[french]{babel}
\usepackage{lmodern}
\usepackage{geometry}
\geometry{margin=1.8cm}
\usepackage{amsmath,amssymb}
\usepackage[most]{tcolorbox}
\usepackage{xcolor}
\usepackage{enumitem}
\usepackage{fancyhdr}
\usepackage{array,tabularx,longtable}
\usepackage{multicol}
\usepackage{tikz}
\usepackage{hyperref}
\usepackage{listings}

%====================================================
% COULEURS
%====================================================
\definecolor{bleu}{RGB}{25,70,140}
\definecolor{vert}{RGB}{30,120,80}
\definecolor{rouge}{RGB}{170,45,45}
\definecolor{orange}{RGB}{210,120,30}
\definecolor{violet}{RGB}{100,70,160}
\definecolor{gris}{RGB}{245,247,250}

%====================================================
% HYPERREF
%====================================================
\hypersetup{
  colorlinks=true,
  linkcolor=bleu,
  urlcolor=bleu
}

%====================================================
% EN-TÊTE
%====================================================
\pagestyle{fancy}
\fancyhf{}
\lhead{\textbf{Programme intensif Brevet 2026}}
\rhead{\textbf{Maths -- Physique-Chimie}}
\cfoot{\thepage}

\setlength{\parindent}{0pt}
\setlength{\parskip}{0.45em}

%====================================================
% ENVIRONNEMENTS
%====================================================
\tcbset{
  enhanced,
  breakable,
  sharp corners,
  boxrule=0.8pt,
  colback=white,
  fonttitle=\bfseries
}

\newtcolorbox{objectif}{
  colframe=bleu,
  colback=blue!3!white,
  title=Objectif
}

\newtcolorbox{methode}{
  colframe=violet,
  colback=violet!4!white,
  title=Méthode
}

\newtcolorbox{retenir}{
  colframe=vert,
  colback=green!3!white,
  title=À retenir
}

\newtcolorbox{attention}{
  colframe=rouge,
  colback=red!3!white,
  title=Attention / erreur classique
}

\newtcolorbox{exemple}{
  colframe=orange,
  colback=orange!4!white,
  title=Exemple corrigé
}

\newcounter{exercice}
\newtcolorbox{exercice}[1][]{
  colframe=black!65,
  colback=gris,
  title={Exercice \refstepcounter{exercice}\theexercice\quad #1}
}

\newcounter{correction}
\newtcolorbox{correction}[1][]{
  colframe=vert!70!black,
  colback=green!2!white,
  title={Correction \refstepcounter{correction}\thecorrection\quad #1}
}

\newcommand{\encadre}[1]{\boxed{\displaystyle #1}}
\newcommand{\kmh}{\,\mathrm{km/h}}
\newcommand{\ms}{\,\mathrm{m/s}}
\newcommand{\cm}{\,\mathrm{cm}}
\newcommand{\m}{\,\mathrm{m}}
\newcommand{\euro}{\,\text{euros}}

%====================================================
% DOCUMENT
%====================================================
\begin{document}
%====================================================
% PAGE DE GARDE
%====================================================
\begin{titlepage}
\centering
\vspace*{1.1cm}

{\Huge \textbf{Programme intensif de révision}}\\[0.3cm]
{\Huge \textbf{Brevet 2026}}\\[0.8cm]

{\LARGE \textbf{Mathématiques et Physique-Chimie}}\\[1cm]

\begin{tcolorbox}[colframe=bleu,colback=blue!3!white,width=0.92\textwidth]
\centering
\Large
Objectif : préparer sérieusement le brevet avec des annales cliquables, un planning jusqu'à l'épreuve, des fiches de méthode, des exercices difficiles et des sujets longs en plusieurs parties.
\end{tcolorbox}

\vspace{0.8cm}

\begin{tcolorbox}[colframe=rouge,colback=red!3!white,width=0.92\textwidth,title=Échéances principales]
\Large
Épreuve de sciences : \textbf{lundi 29 juin 2026}.\\
Épreuve de mathématiques : \textbf{mardi 30 juin 2026, 9h--11h}.
\end{tcolorbox}

\vfill

{\Large Niveau troisième générale}\\[0.2cm]
{\Large Livret professeur -- élève}\\[0.5cm]

\end{titlepage}

\tableofcontents
\newpage
\section{Mode d'emploi du livret}

\begin{objectif}
Ce document sert à organiser une révision complète jusqu'au brevet. Il contient :
\begin{itemize}
  \item des liens cliquables vers des sujets de brevet des années précédentes ;
  \item un planning jour par jour jusqu'au brevet ;
  \item des fiches de méthode en mathématiques et en physique-chimie ;
  \item des sujets inventés longs, en plusieurs parties ;
  \item des corrections détaillées ou semi-détaillées.
\end{itemize}
\end{objectif}

\begin{methode}
\textbf{Méthode de travail}
\begin{enumerate}
  \item Faire les exercices sans regarder la correction.
  \item Se mettre en temps limité pour les annales.
  \item Corriger en distinguant les erreurs de cours, de méthode, de calcul et de rédaction.
  \item Refaire les exercices ratés deux jours plus tard.
  \item Tenir une fiche d'erreurs personnelle.
\end{enumerate}
\end{methode}

\begin{attention}
Le brevet ne récompense pas uniquement le résultat final. Les étapes, les justifications, les unités et les phrases de conclusion comptent beaucoup.
\end{attention}

\section{Annales officielles et sujets de brevet cliquables}

\begin{objectif}
Les liens suivants permettent de travailler sur de vrais sujets de brevet de mathématiques des années précédentes. Ils sont cliquables dans le PDF.
\end{objectif}

\begin{retenir}
\textbf{Banques générales d'annales :}
\begin{itemize}
  \item \href{https://www.apmep.fr/Annales-du-Brevet-des-colleges}{APMEP -- Annales du brevet des collèges}.
  \item \href{https://www.apmep.fr/Annales-du-Brevet-cycle-4}{APMEP -- Annales du brevet cycle 4}.
  \item \href{https://groupe-reussite.fr/ressources/annales-brevet-maths/}{Groupe Réussite -- Annales de brevet de mathématiques}.
  \item \href{https://www.lumni.fr/article/maths-sujets-et-corriges-serie-generale-brevet-2025}{Lumni -- Sujet et corrigé de mathématiques du brevet 2025}.
  \item \href{https://www.letudiant.fr/college/brevet/sujets-et-corriges-du-brevet-de-mathematiques-2025-dnb.html}{L'Étudiant -- Sujet et corrigé du brevet de mathématiques 2025}.
\end{itemize}
\end{retenir}

\begin{retenir}
\textbf{Annales APMEP par année :}
\begin{itemize}
  \item \href{https://www.apmep.fr/Brevet-2025}{Brevet 2025 et sujets zéro 2026}.
  \item \href{https://www.apmep.fr/Brevet-2024}{Brevet 2024}.
  \item \href{https://www.apmep.fr/Brevet-2023}{Brevet 2023}.
  \item \href{https://www.apmep.fr/Brevet-2022}{Brevet 2022}.
  \item \href{https://www.apmep.fr/Brevet-2021}{Brevet 2021}.
  \item \href{https://www.apmep.fr/Brevet-2020}{Brevet 2020}.
  \item \href{https://www.apmep.fr/Brevet-2019}{Brevet 2019}.
  \item \href{https://www.apmep.fr/Brevet-2018}{Brevet 2018}.
\end{itemize}
\end{retenir}

\begin{retenir}
\textbf{Intégrales PDF utiles :}
\begin{itemize}
  \item \href{https://www.apmep.fr/IMG/pdf/Annee_2025_Brevet.pdf}{APMEP -- Intégrale brevet 2025}.
  \item \href{https://www.apmep.fr/IMG/pdf/Brevet_annee_2024_DV.pdf}{APMEP -- Intégrale brevet 2024}.
  \item \href{https://www.apmep.fr/IMG/pdf/Brevet_Annee_2023_DV_2.pdf}{APMEP -- Intégrale brevet 2023}.
\end{itemize}
\end{retenir}

\begin{methode}
\textbf{Comment travailler une annale}
\begin{enumerate}
  \item Faire le sujet en 2 h, sans correction.
  \item Garder 10 minutes de relecture.
  \item Corriger avec une couleur différente.
  \item Refaire immédiatement les questions où la méthode était connue mais mal appliquée.
  \item Refaire 48 h plus tard les questions vraiment ratées.
\end{enumerate}
\end{methode}

\newpage
\section{Planning intensif jusqu'au brevet}

\begin{objectif}
Le planning suivant part du vendredi 29 mai 2026 et va jusqu'à l'épreuve de mathématiques du mardi 30 juin 2026. Il alterne cours, automatismes, annales, correction active et sujets inventés difficiles.
\end{objectif}

\renewcommand{\arraystretch}{1.35}
\begin{longtable}{|p{2.3cm}|p{3.2cm}|p{4.7cm}|p{4.7cm}|p{0.7cm}|}
\hline
\textbf{Date} & \textbf{Objectif} & \textbf{Mathématiques} & \textbf{Physique-Chimie} & \textbf{Fait} \\
\hline
Ven. 29 mai & Diagnostic initial & Calculs courts : fractions, puissances, priorités. Puis M-LONG-01, partie A. & Conversions : km/h, m/s, L, mL, g, kg. PC-LONG-01, partie A. & $\square$ \\
\hline
Sam. 30 mai & Calcul numérique & Fractions, puissances de 10, notation scientifique, racines carrées. M-LONG-01, partie B. & Mouvement, vitesse moyenne, distance, durée. PC-LONG-01, partie B. & $\square$ \\
\hline
Dim. 31 mai & Correction active & Refaire toutes les erreurs du diagnostic. Faire une fiche : calculs dangereux. & Refaire les conversions et formules ratées. & $\square$ \\
\hline
Lun. 1 juin & Calcul littéral & Développer, réduire, factoriser. M-LONG-01, partie C. & Électricité : tension, intensité, résistance, loi d'Ohm. & $\square$ \\
\hline
Mar. 2 juin & Équations & Équations du premier degré, équations produit, mise en équation. Faire M-LONG-01 en entier. & Puissance et énergie électrique. & $\square$ \\
\hline
Mer. 3 juin & Fonctions & Images, antécédents, lecture graphique, fonctions affines. M-LONG-02, partie A. & Signaux : son, lumière, vitesse de propagation. & $\square$ \\
\hline
Jeu. 4 juin & Proportionnalité & Pourcentages, échelles, coefficients multiplicateurs, vitesse. M-LONG-02, partie B. & Masse volumique, densité, identification d'un matériau. & $\square$ \\
\hline
Ven. 5 juin & Géométrie plane & Pythagore, réciproque de Pythagore, trigonométrie. M-LONG-02, partie C. & Atomes, ions, molécules. & $\square$ \\
\hline
Sam. 6 juin & Thalès & Thalès, réciproque, triangles semblables, agrandissement-réduction. M-LONG-02 en entier. & Transformations chimiques et équations équilibrées. & $\square$ \\
\hline
Dim. 7 juin & Annale 2025 & Faire un sujet sur \href{https://www.apmep.fr/Brevet-2025}{APMEP -- Brevet 2025}. Durée : 2 h. & Relecture fiches sciences. & $\square$ \\
\hline
Lun. 8 juin & Correction annale & Corriger l'annale 2025. Refaire deux exercices ratés sans correction. & Refaire les questions d'unités ratées. & $\square$ \\
\hline
Mar. 9 juin & Statistiques & Moyenne, médiane, étendue, fréquences. M-LONG-03, partie A. & Énergie : joule, kilowattheure, coût. & $\square$ \\
\hline
Mer. 10 juin & Probabilités & Probabilités simples, événement contraire, tableaux à deux entrées. M-LONG-03, partie B. & Loi d'Ohm et circuits simples. & $\square$ \\
\hline
Jeu. 11 juin & Algorithmique & Programmes Scratch/Python, boucles, variables, seuils. M-LONG-03, partie C. & Signaux : écho, distance, vitesse du son. & $\square$ \\
\hline
Ven. 12 juin & Annale 2024 & Faire un sujet sur \href{https://www.apmep.fr/Brevet-2024}{APMEP -- Brevet 2024}. Durée : 2 h. & Mini-sujet sciences inventé. & $\square$ \\
\hline
Sam. 13 juin & Correction annale & Corriger l'annale 2024. Remplir le carnet d'erreurs. & Refaire les exercices sciences ratés. & $\square$ \\
\hline
Dim. 14 juin & Repos actif & 45 min : formules essentielles et erreurs classiques. & 30 min : unités et équations chimiques. & $\square$ \\
\hline
Lun. 15 juin & Sujet inventé difficile & Faire M-LONG-03 en entier en temps limité. & Faire PC-LONG-01 en entier. & $\square$ \\
\hline
Mar. 16 juin & Géométrie espace & Volumes, sphère, cylindre, cône, sections, agrandissement de volumes. M-LONG-04, partie A. & Masse volumique et énergie. & $\square$ \\
\hline
Mer. 17 juin & Annale 2023 & Faire un sujet sur \href{https://www.apmep.fr/Brevet-2023}{APMEP -- Brevet 2023}. Durée : 2 h. & Relecture chimie. & $\square$ \\
\hline
Jeu. 18 juin & Correction annale & Correction active de l'annale 2023. Refaire les questions perdues. & Électricité + énergie. & $\square$ \\
\hline
Ven. 19 juin & Sujet transversal & Faire M-LONG-04, parties B et C. & Faire PC-LONG-02, partie A. & $\square$ \\
\hline
Sam. 20 juin & Brevet blanc maison & Faire BB-MATHS-01. Durée : 2 h. & Faire BB-PC-01. Durée : 1 h. & $\square$ \\
\hline
Dim. 21 juin & Correction brevet blanc & Corriger le brevet blanc. Refaire les exercices ratés. & Corriger le sujet sciences. & $\square$ \\
\hline
Lun. 22 juin & Fonctions et équations & Refaire tous les exercices de fonctions et équations ratés. & Revoir loi d'Ohm et énergie. & $\square$ \\
\hline
Mar. 23 juin & Problèmes complexes & Faire M-LONG-05, sujet transversal sur un voyage scolaire. & Faire PC-LONG-02, partie B. & $\square$ \\
\hline
Mer. 24 juin & Annale 2022 & Faire un sujet sur \href{https://www.apmep.fr/Brevet-2022}{APMEP -- Brevet 2022}. Durée : 2 h. & Mini-sujet sciences. & $\square$ \\
\hline
Jeu. 25 juin & Correction finale & Corriger l'annale 2022. Ne plus découvrir de nouveau chapitre. & Relecture complète sciences. & $\square$ \\
\hline
Ven. 26 juin & Épreuve de français & Maths léger : 30 min de formules seulement. & Sciences léger : unités et schémas. & $\square$ \\
\hline
Sam. 27 juin & Dernier sujet & Faire un dernier sujet représentatif sur \href{https://www.lumni.fr/article/maths-sujets-et-corriges-serie-generale-brevet-2025}{Lumni -- Brevet maths 2025}. & Revoir physique-chimie. & $\square$ \\
\hline
Dim. 28 juin & Repos intelligent & Refaire uniquement les erreurs déjà corrigées. Pas de nouveau sujet difficile. & Dernière fiche sciences. & $\square$ \\
\hline
Lun. 29 juin & Épreuve de sciences & Le soir : 45 min de formules maths, pas plus. & Épreuve de sciences. & $\square$ \\
\hline
Mar. 30 juin & Épreuve de maths & Brevet de mathématiques : commencer par le plus rentable, rédiger, relire. & Repos après l'épreuve. & $\square$ \\
\hline
\end{longtable}
\newpage
\section{Fiches méthodes express}

\subsection{Calcul numérique}
\begin{retenir}
\textbf{Priorités opératoires :} parenthèses, puissances, multiplications et divisions, additions et soustractions.
\[
\frac{a}{b}+\frac{c}{b}=\frac{a+c}{b},
\qquad
\frac{a}{b}\times\frac{c}{d}=\frac{ac}{bd},
\qquad
\frac{a}{b}\div\frac{c}{d}=\frac{a}{b}\times\frac{d}{c}.
\]
\[
a^m\times a^n=a^{m+n},
\qquad
\frac{a^m}{a^n}=a^{m-n},
\qquad
(a^m)^n=a^{mn}.
\]
\end{retenir}

\begin{attention}
\[
(-4)^2=16
\quad \text{mais} \quad
-4^2=-16.
\]
Les parenthèses changent le nombre qui est mis au carré.
\end{attention}

\subsection{Calcul littéral et équations}
\begin{retenir}
\[
k(a+b)=ka+kb,
\qquad
(a+b)(c+d)=ac+ad+bc+bd.
\]
\[
(a+b)^2=a^2+2ab+b^2,
\qquad
(a-b)^2=a^2-2ab+b^2,
\qquad
(a+b)(a-b)=a^2-b^2.
\]
\end{retenir}

\begin{methode}
Pour résoudre une équation : développer si nécessaire, réduire, regrouper les termes avec l'inconnue, isoler l'inconnue, vérifier.
\end{methode}

\subsection{Fonctions}
\begin{retenir}
Une fonction linéaire est de la forme $f(x)=ax$. Elle modélise une proportionnalité. Une fonction affine est de la forme $f(x)=ax+b$.
\end{retenir}

\subsection{Géométrie}
\begin{retenir}
\[
\text{Pythagore :}\quad c^2=a^2+b^2.
\]
\[
\sin(\widehat{A})=\frac{\text{opposé}}{\text{hypoténuse}},
\qquad
\cos(\widehat{A})=\frac{\text{adjacent}}{\text{hypoténuse}},
\qquad
\tan(\widehat{A})=\frac{\text{opposé}}{\text{adjacent}}.
\]
\[
\text{Aire triangle}=\frac{b\times h}{2},
\qquad
\text{Volume cylindre}=\pi r^2h.
\]
\end{retenir}

\subsection{Physique-Chimie}
\begin{retenir}
\[
v=\frac{d}{t},
\qquad
d=v\times t,
\qquad
E=P\times t,
\qquad
U=RI,
\qquad
\rho=\frac{m}{V}.
\]
\[
1\ \mathrm{h}=3600\ \mathrm{s},
\qquad
1\ \mathrm{km}=1000\ \mathrm{m},
\qquad
v(\mathrm{m/s})=\frac{v(\mathrm{km/h})}{3{,}6}.
\]
\end{retenir}
\newpage
\section{Sujets inventés longs et difficiles de mathématiques}

\begin{attention}
Les sujets suivants sont inventés. Ils sont plus longs et parfois plus difficiles qu'un exercice classique de brevet. Ils servent à préparer les bons élèves, à travailler la prise d'initiative et à consolider toutes les méthodes du programme de troisième.
\end{attention}
\subsection{M-LONG-01 -- Calculs, programme de calcul et équations}

\begin{objectif}
Durée conseillée : 1 h à 1 h 45.\\
Chapitres : fractions, puissances, calcul littéral, équations, équations produit.
\end{objectif}
\begin{exercice}[M-LONG-01 -- Partie A : calcul numérique]
On considère les nombres :
\[
A=7-3(2-5)^2+\frac{5}{6}\div\frac{10}{9},
\qquad
B=\frac{2^5\times 2^{-3}}{4}+\frac{3}{5}-\frac{7}{10},
\]
\[
C=\left(\frac{3}{4}-\frac{5}{6}\right)^2.
\]
\begin{enumerate}
  \item Calculer $(2-5)^2$.
  \item Calculer $\dfrac{5}{6}\div\dfrac{10}{9}$.
  \item En déduire la valeur exacte de $A$.
  \item Simplifier $2^5\times2^{-3}$.
  \item Calculer $B$ sous forme d'une fraction irréductible.
  \item Calculer $C$ sous forme d'une fraction irréductible.
  \item Ranger $A$, $B$ et $C$ dans l'ordre croissant.
\end{enumerate}
\end{exercice}

\begin{exercice}[M-LONG-01 -- Partie B : programme de calcul]
On considère le programme de calcul suivant : choisir un nombre ; ajouter $4$ ; multiplier le résultat par le nombre choisi au départ ; soustraire le carré du nombre choisi ; ajouter $7$.
\begin{enumerate}
  \item Appliquer ce programme au nombre $3$.
  \item Appliquer ce programme au nombre $-5$.
  \item On note $x$ le nombre choisi au départ. Montrer que le résultat final est $4x+7$.
  \item Quel nombre faut-il choisir pour obtenir $51$ ?
  \item Peut-on obtenir $-20$ ? Justifier.
\end{enumerate}
\end{exercice}

\begin{exercice}[M-LONG-01 -- Partie C : expression factorisée]
On considère :
\[
E=(x+4)^2-(x+4)(2x-1).
\]
\begin{enumerate}
  \item Développer et réduire $E$.
  \item Factoriser $E$ sans développer.
  \item Résoudre l'équation $E=0$.
  \item Résoudre l'équation $E=9$ en utilisant la forme développée.
  \item Vérifier une solution obtenue à la question précédente.
\end{enumerate}
\end{exercice}
\newpage
\subsection{M-LONG-02 -- Fonctions, tarifs et proportionnalité}

\begin{objectif}
Durée conseillée : 1 h à 1 h 45.\\
Chapitres : fonctions affines, proportionnalité, pourcentages, inéquations.
\end{objectif}
\begin{exercice}[M-LONG-02 -- Partie A : deux tarifs]
Une ville propose deux tarifs pour louer un vélo électrique : tarif A, $0{,}30$ euro par minute ; tarif B, abonnement de $5$ euros puis $0{,}12$ euro par minute. On note $x$ le nombre de minutes.
\begin{enumerate}
  \item Écrire $A(x)$ et $B(x)$.
  \item Calculer $A(20)$, $B(20)$, $A(45)$ et $B(45)$.
  \item Résoudre $B(x)<A(x)$.
  \item À partir de quelle durée entière le tarif B devient-il plus avantageux ?
\end{enumerate}
\end{exercice}

\begin{exercice}[M-LONG-02 -- Partie B : lecture et interprétation]
On considère $A(x)=0{,}30x$ et $B(x)=5+0{,}12x$.
\begin{enumerate}
  \item Dire si $A$ est linéaire ou affine.
  \item Dire si $B$ est linéaire ou affine.
  \item Que représentent $0{,}30$, $0{,}12$ et $5$ dans le contexte ?
  \item Résoudre $A(x)=B(x)$.
  \item Interpréter graphiquement le résultat.
\end{enumerate}
\end{exercice}

\begin{exercice}[M-LONG-02 -- Partie C : augmentation de prix]
L'année suivante, la ville augmente le tarif A de $10\%$, l'abonnement du tarif B de $20\%$ et le prix par minute du tarif B de $5\%$.
\begin{enumerate}
  \item Écrire les nouvelles fonctions $A_2(x)$ et $B_2(x)$.
  \item Comparer les deux nouveaux tarifs pour $40$ minutes.
  \item Résoudre $B_2(x)<A_2(x)$.
  \item Expliquer l'effet de ces augmentations sur le choix du tarif.
\end{enumerate}
\end{exercice}
\newpage
\subsection{M-LONG-03 -- Statistiques, probabilités et algorithmique}

\begin{objectif}
Durée conseillée : 1 h à 1 h 45.\\
Chapitres : moyenne, médiane, étendue, probabilités, algorithmes.
\end{objectif}
\begin{exercice}[M-LONG-03 -- Partie A : statistiques]
Une application de révision enregistre le nombre de questions réussies par 15 élèves :
\[
12,15,9,18,20,14,16,10,11,13,19,17,15,8,14.
\]
\begin{enumerate}
  \item Calculer l'étendue.
  \item Calculer la moyenne.
  \item Ranger les valeurs dans l'ordre croissant.
  \item Déterminer la médiane.
  \item Calculer la fréquence des élèves ayant réussi au moins 15 questions.
\end{enumerate}
\end{exercice}

\begin{exercice}[M-LONG-03 -- Partie B : probabilités]
Une question est choisie parmi 5 questions de calcul, 4 de géométrie, 3 de probabilités et 2 de physique.
\begin{enumerate}
  \item Combien y a-t-il de questions au total ?
  \item Quelle est la probabilité de tomber sur une question de géométrie ?
  \item Quelle est la probabilité de ne pas tomber sur une question de physique ?
  \item On tire une question, puis une deuxième avec remise. Quelle est la probabilité d'obtenir deux questions de calcul ?
  \item Quelle est la probabilité d'obtenir au moins une question de physique en deux tirages ?
\end{enumerate}
\end{exercice}

\begin{exercice}[M-LONG-03 -- Partie C : algorithme de seuil]
Un élève gagne $18$ points le premier jour. Chaque jour suivant, il gagne $4$ points de plus que la veille. On cherche le premier jour où le total dépasse $500$ points.
\begin{enumerate}
  \item Compléter un tableau jusqu'au jour 5.
  \item Écrire une méthode de calcul progressive.
  \item Expliquer le rôle des variables \texttt{total}, \texttt{jour} et \texttt{points}.
  \item Par essais successifs, déterminer le premier jour où le total dépasse $500$ points.
\end{enumerate}
\end{exercice}
\newpage
\subsection{M-LONG-04 -- Géométrie, volumes et optimisation}

\begin{objectif}
Durée conseillée : 1 h à 1 h 45.\\
Chapitres : Pythagore, trigonométrie, volumes, agrandissements.
\end{objectif}
\begin{exercice}[M-LONG-04 -- Partie A : étude de situation]
On étudie une situation contextualisée intitulée \og Géométrie, volumes et optimisation \fg. Les données sont les suivantes : une grandeur de départ vaut $120$, elle évolue selon un coefficient multiplicateur $1{,}08$, puis on lui retranche $15$ unités fixes. Une autre grandeur est donnée par la fonction $f(x)=3x+24$.
\begin{enumerate}
  \item Expliquer pourquoi une augmentation de $8\%$ correspond à une multiplication par $1{,}08$.
  \item Calculer la valeur obtenue après augmentation à partir de $120$.
  \item Retirer ensuite $15$ unités et donner le résultat final.
  \item Calculer $f(0)$, $f(12)$ et $f(25)$.
  \item Déterminer l'antécédent de $99$ par la fonction $f$.
  \item Interpréter ce résultat dans le contexte du sujet.
\end{enumerate}
\end{exercice}

\begin{exercice}[M-LONG-04 -- Partie B : calcul littéral et équation]
On considère l'expression :
\[
E=(2x-5)^2-(x+3)(x-3).
\]
\begin{enumerate}
  \item Développer $(2x-5)^2$.
  \item Développer $(x+3)(x-3)$.
  \item En déduire une forme développée et réduite de $E$.
  \item Calculer $E$ pour $x=4$.
  \item Résoudre l'équation $E=25$.
  \item Vérifier une solution obtenue.
\end{enumerate}
\end{exercice}

\begin{exercice}[M-LONG-04 -- Partie C : géométrie et grandeurs]
Un triangle rectangle est utilisé dans la modélisation. Ses côtés de l'angle droit mesurent $9\ \mathrm{cm}$ et $12\ \mathrm{cm}$.
\begin{enumerate}
  \item Calculer la longueur de l'hypoténuse.
  \item Calculer l'aire du triangle.
  \item Calculer $\cos(\widehat{B})$ pour l'angle adjacent au côté de longueur $9\ \mathrm{cm}$.
  \item En déduire une valeur approchée de l'angle.
  \item On agrandit la figure avec un coefficient $1{,}4$. Calculer les nouvelles longueurs.
  \item Par combien l'aire est-elle multipliée ?
\end{enumerate}
\end{exercice}

\begin{exercice}[M-LONG-04 -- Partie D : synthèse et prise d'initiative]
Une décision doit être prise à partir des résultats des parties précédentes.
\begin{enumerate}
  \item Rédiger une conclusion de cinq lignes qui utilise au moins deux résultats numériques obtenus.
  \item Vérifier la cohérence des unités.
  \item Proposer une amélioration du modèle.
  \item Donner une limite du raisonnement.
  \item Écrire une phrase finale comme dans une copie de brevet.
\end{enumerate}
\end{exercice}
\newpage
\subsection{M-LONG-05 -- Voyage scolaire}

\begin{objectif}
Durée conseillée : 1 h à 1 h 45.\\
Chapitres : fonctions, vitesse, équations, probabilités.
\end{objectif}
\begin{exercice}[M-LONG-05 -- Partie A : étude de situation]
On étudie une situation contextualisée intitulée \og Voyage scolaire \fg. Les données sont les suivantes : une grandeur de départ vaut $120$, elle évolue selon un coefficient multiplicateur $1{,}08$, puis on lui retranche $15$ unités fixes. Une autre grandeur est donnée par la fonction $f(x)=3x+24$.
\begin{enumerate}
  \item Expliquer pourquoi une augmentation de $8\%$ correspond à une multiplication par $1{,}08$.
  \item Calculer la valeur obtenue après augmentation à partir de $120$.
  \item Retirer ensuite $15$ unités et donner le résultat final.
  \item Calculer $f(0)$, $f(12)$ et $f(25)$.
  \item Déterminer l'antécédent de $99$ par la fonction $f$.
  \item Interpréter ce résultat dans le contexte du sujet.
\end{enumerate}
\end{exercice}

\begin{exercice}[M-LONG-05 -- Partie B : calcul littéral et équation]
On considère l'expression :
\[
E=(2x-5)^2-(x+3)(x-3).
\]
\begin{enumerate}
  \item Développer $(2x-5)^2$.
  \item Développer $(x+3)(x-3)$.
  \item En déduire une forme développée et réduite de $E$.
  \item Calculer $E$ pour $x=4$.
  \item Résoudre l'équation $E=25$.
  \item Vérifier une solution obtenue.
\end{enumerate}
\end{exercice}

\begin{exercice}[M-LONG-05 -- Partie C : géométrie et grandeurs]
Un triangle rectangle est utilisé dans la modélisation. Ses côtés de l'angle droit mesurent $9\ \mathrm{cm}$ et $12\ \mathrm{cm}$.
\begin{enumerate}
  \item Calculer la longueur de l'hypoténuse.
  \item Calculer l'aire du triangle.
  \item Calculer $\cos(\widehat{B})$ pour l'angle adjacent au côté de longueur $9\ \mathrm{cm}$.
  \item En déduire une valeur approchée de l'angle.
  \item On agrandit la figure avec un coefficient $1{,}4$. Calculer les nouvelles longueurs.
  \item Par combien l'aire est-elle multipliée ?
\end{enumerate}
\end{exercice}

\begin{exercice}[M-LONG-05 -- Partie D : synthèse et prise d'initiative]
Une décision doit être prise à partir des résultats des parties précédentes.
\begin{enumerate}
  \item Rédiger une conclusion de cinq lignes qui utilise au moins deux résultats numériques obtenus.
  \item Vérifier la cohérence des unités.
  \item Proposer une amélioration du modèle.
  \item Donner une limite du raisonnement.
  \item Écrire une phrase finale comme dans une copie de brevet.
\end{enumerate}
\end{exercice}
\newpage
\subsection{M-LONG-06 -- Piscine municipale et consommation d'eau}

\begin{objectif}
Durée conseillée : 1 h à 1 h 45.\\
Chapitres : volumes, pourcentages, fonctions, vitesse de remplissage.
\end{objectif}
\begin{exercice}[M-LONG-06 -- Partie A : étude de situation]
On étudie une situation contextualisée intitulée \og Piscine municipale et consommation d'eau \fg. Les données sont les suivantes : une grandeur de départ vaut $120$, elle évolue selon un coefficient multiplicateur $1{,}08$, puis on lui retranche $15$ unités fixes. Une autre grandeur est donnée par la fonction $f(x)=3x+24$.
\begin{enumerate}
  \item Expliquer pourquoi une augmentation de $8\%$ correspond à une multiplication par $1{,}08$.
  \item Calculer la valeur obtenue après augmentation à partir de $120$.
  \item Retirer ensuite $15$ unités et donner le résultat final.
  \item Calculer $f(0)$, $f(12)$ et $f(25)$.
  \item Déterminer l'antécédent de $99$ par la fonction $f$.
  \item Interpréter ce résultat dans le contexte du sujet.
\end{enumerate}
\end{exercice}

\begin{exercice}[M-LONG-06 -- Partie B : calcul littéral et équation]
On considère l'expression :
\[
E=(2x-5)^2-(x+3)(x-3).
\]
\begin{enumerate}
  \item Développer $(2x-5)^2$.
  \item Développer $(x+3)(x-3)$.
  \item En déduire une forme développée et réduite de $E$.
  \item Calculer $E$ pour $x=4$.
  \item Résoudre l'équation $E=25$.
  \item Vérifier une solution obtenue.
\end{enumerate}
\end{exercice}

\begin{exercice}[M-LONG-06 -- Partie C : géométrie et grandeurs]
Un triangle rectangle est utilisé dans la modélisation. Ses côtés de l'angle droit mesurent $9\ \mathrm{cm}$ et $12\ \mathrm{cm}$.
\begin{enumerate}
  \item Calculer la longueur de l'hypoténuse.
  \item Calculer l'aire du triangle.
  \item Calculer $\cos(\widehat{B})$ pour l'angle adjacent au côté de longueur $9\ \mathrm{cm}$.
  \item En déduire une valeur approchée de l'angle.
  \item On agrandit la figure avec un coefficient $1{,}4$. Calculer les nouvelles longueurs.
  \item Par combien l'aire est-elle multipliée ?
\end{enumerate}
\end{exercice}

\begin{exercice}[M-LONG-06 -- Partie D : synthèse et prise d'initiative]
Une décision doit être prise à partir des résultats des parties précédentes.
\begin{enumerate}
  \item Rédiger une conclusion de cinq lignes qui utilise au moins deux résultats numériques obtenus.
  \item Vérifier la cohérence des unités.
  \item Proposer une amélioration du modèle.
  \item Donner une limite du raisonnement.
  \item Écrire une phrase finale comme dans une copie de brevet.
\end{enumerate}
\end{exercice}
\newpage
\subsection{M-LONG-07 -- Parc solaire du collège}

\begin{objectif}
Durée conseillée : 1 h à 1 h 45.\\
Chapitres : proportionnalité, statistiques, énergie, fonctions.
\end{objectif}
\begin{exercice}[M-LONG-07 -- Partie A : étude de situation]
On étudie une situation contextualisée intitulée \og Parc solaire du collège \fg. Les données sont les suivantes : une grandeur de départ vaut $120$, elle évolue selon un coefficient multiplicateur $1{,}08$, puis on lui retranche $15$ unités fixes. Une autre grandeur est donnée par la fonction $f(x)=3x+24$.
\begin{enumerate}
  \item Expliquer pourquoi une augmentation de $8\%$ correspond à une multiplication par $1{,}08$.
  \item Calculer la valeur obtenue après augmentation à partir de $120$.
  \item Retirer ensuite $15$ unités et donner le résultat final.
  \item Calculer $f(0)$, $f(12)$ et $f(25)$.
  \item Déterminer l'antécédent de $99$ par la fonction $f$.
  \item Interpréter ce résultat dans le contexte du sujet.
\end{enumerate}
\end{exercice}

\begin{exercice}[M-LONG-07 -- Partie B : calcul littéral et équation]
On considère l'expression :
\[
E=(2x-5)^2-(x+3)(x-3).
\]
\begin{enumerate}
  \item Développer $(2x-5)^2$.
  \item Développer $(x+3)(x-3)$.
  \item En déduire une forme développée et réduite de $E$.
  \item Calculer $E$ pour $x=4$.
  \item Résoudre l'équation $E=25$.
  \item Vérifier une solution obtenue.
\end{enumerate}
\end{exercice}

\begin{exercice}[M-LONG-07 -- Partie C : géométrie et grandeurs]
Un triangle rectangle est utilisé dans la modélisation. Ses côtés de l'angle droit mesurent $9\ \mathrm{cm}$ et $12\ \mathrm{cm}$.
\begin{enumerate}
  \item Calculer la longueur de l'hypoténuse.
  \item Calculer l'aire du triangle.
  \item Calculer $\cos(\widehat{B})$ pour l'angle adjacent au côté de longueur $9\ \mathrm{cm}$.
  \item En déduire une valeur approchée de l'angle.
  \item On agrandit la figure avec un coefficient $1{,}4$. Calculer les nouvelles longueurs.
  \item Par combien l'aire est-elle multipliée ?
\end{enumerate}
\end{exercice}

\begin{exercice}[M-LONG-07 -- Partie D : synthèse et prise d'initiative]
Une décision doit être prise à partir des résultats des parties précédentes.
\begin{enumerate}
  \item Rédiger une conclusion de cinq lignes qui utilise au moins deux résultats numériques obtenus.
  \item Vérifier la cohérence des unités.
  \item Proposer une amélioration du modèle.
  \item Donner une limite du raisonnement.
  \item Écrire une phrase finale comme dans une copie de brevet.
\end{enumerate}
\end{exercice}
\newpage
\subsection{M-LONG-08 -- Course d'orientation}

\begin{objectif}
Durée conseillée : 1 h à 1 h 45.\\
Chapitres : géométrie, échelle, vitesse, trigonométrie.
\end{objectif}
\begin{exercice}[M-LONG-08 -- Partie A : étude de situation]
On étudie une situation contextualisée intitulée \og Course d'orientation \fg. Les données sont les suivantes : une grandeur de départ vaut $120$, elle évolue selon un coefficient multiplicateur $1{,}08$, puis on lui retranche $15$ unités fixes. Une autre grandeur est donnée par la fonction $f(x)=3x+24$.
\begin{enumerate}
  \item Expliquer pourquoi une augmentation de $8\%$ correspond à une multiplication par $1{,}08$.
  \item Calculer la valeur obtenue après augmentation à partir de $120$.
  \item Retirer ensuite $15$ unités et donner le résultat final.
  \item Calculer $f(0)$, $f(12)$ et $f(25)$.
  \item Déterminer l'antécédent de $99$ par la fonction $f$.
  \item Interpréter ce résultat dans le contexte du sujet.
\end{enumerate}
\end{exercice}

\begin{exercice}[M-LONG-08 -- Partie B : calcul littéral et équation]
On considère l'expression :
\[
E=(2x-5)^2-(x+3)(x-3).
\]
\begin{enumerate}
  \item Développer $(2x-5)^2$.
  \item Développer $(x+3)(x-3)$.
  \item En déduire une forme développée et réduite de $E$.
  \item Calculer $E$ pour $x=4$.
  \item Résoudre l'équation $E=25$.
  \item Vérifier une solution obtenue.
\end{enumerate}
\end{exercice}

\begin{exercice}[M-LONG-08 -- Partie C : géométrie et grandeurs]
Un triangle rectangle est utilisé dans la modélisation. Ses côtés de l'angle droit mesurent $9\ \mathrm{cm}$ et $12\ \mathrm{cm}$.
\begin{enumerate}
  \item Calculer la longueur de l'hypoténuse.
  \item Calculer l'aire du triangle.
  \item Calculer $\cos(\widehat{B})$ pour l'angle adjacent au côté de longueur $9\ \mathrm{cm}$.
  \item En déduire une valeur approchée de l'angle.
  \item On agrandit la figure avec un coefficient $1{,}4$. Calculer les nouvelles longueurs.
  \item Par combien l'aire est-elle multipliée ?
\end{enumerate}
\end{exercice}

\begin{exercice}[M-LONG-08 -- Partie D : synthèse et prise d'initiative]
Une décision doit être prise à partir des résultats des parties précédentes.
\begin{enumerate}
  \item Rédiger une conclusion de cinq lignes qui utilise au moins deux résultats numériques obtenus.
  \item Vérifier la cohérence des unités.
  \item Proposer une amélioration du modèle.
  \item Donner une limite du raisonnement.
  \item Écrire une phrase finale comme dans une copie de brevet.
\end{enumerate}
\end{exercice}
\newpage
\subsection{M-LONG-09 -- Application de révision}

\begin{objectif}
Durée conseillée : 1 h à 1 h 45.\\
Chapitres : statistiques, fonctions, algorithme, probabilités.
\end{objectif}
\begin{exercice}[M-LONG-09 -- Partie A : étude de situation]
On étudie une situation contextualisée intitulée \og Application de révision \fg. Les données sont les suivantes : une grandeur de départ vaut $120$, elle évolue selon un coefficient multiplicateur $1{,}08$, puis on lui retranche $15$ unités fixes. Une autre grandeur est donnée par la fonction $f(x)=3x+24$.
\begin{enumerate}
  \item Expliquer pourquoi une augmentation de $8\%$ correspond à une multiplication par $1{,}08$.
  \item Calculer la valeur obtenue après augmentation à partir de $120$.
  \item Retirer ensuite $15$ unités et donner le résultat final.
  \item Calculer $f(0)$, $f(12)$ et $f(25)$.
  \item Déterminer l'antécédent de $99$ par la fonction $f$.
  \item Interpréter ce résultat dans le contexte du sujet.
\end{enumerate}
\end{exercice}

\begin{exercice}[M-LONG-09 -- Partie B : calcul littéral et équation]
On considère l'expression :
\[
E=(2x-5)^2-(x+3)(x-3).
\]
\begin{enumerate}
  \item Développer $(2x-5)^2$.
  \item Développer $(x+3)(x-3)$.
  \item En déduire une forme développée et réduite de $E$.
  \item Calculer $E$ pour $x=4$.
  \item Résoudre l'équation $E=25$.
  \item Vérifier une solution obtenue.
\end{enumerate}
\end{exercice}

\begin{exercice}[M-LONG-09 -- Partie C : géométrie et grandeurs]
Un triangle rectangle est utilisé dans la modélisation. Ses côtés de l'angle droit mesurent $9\ \mathrm{cm}$ et $12\ \mathrm{cm}$.
\begin{enumerate}
  \item Calculer la longueur de l'hypoténuse.
  \item Calculer l'aire du triangle.
  \item Calculer $\cos(\widehat{B})$ pour l'angle adjacent au côté de longueur $9\ \mathrm{cm}$.
  \item En déduire une valeur approchée de l'angle.
  \item On agrandit la figure avec un coefficient $1{,}4$. Calculer les nouvelles longueurs.
  \item Par combien l'aire est-elle multipliée ?
\end{enumerate}
\end{exercice}

\begin{exercice}[M-LONG-09 -- Partie D : synthèse et prise d'initiative]
Une décision doit être prise à partir des résultats des parties précédentes.
\begin{enumerate}
  \item Rédiger une conclusion de cinq lignes qui utilise au moins deux résultats numériques obtenus.
  \item Vérifier la cohérence des unités.
  \item Proposer une amélioration du modèle.
  \item Donner une limite du raisonnement.
  \item Écrire une phrase finale comme dans une copie de brevet.
\end{enumerate}
\end{exercice}
\newpage
\subsection{M-LONG-10 -- Réaménagement d'une cour}

\begin{objectif}
Durée conseillée : 1 h à 1 h 45.\\
Chapitres : aires, périmètres, équations, optimisation.
\end{objectif}
\begin{exercice}[M-LONG-10 -- Partie A : étude de situation]
On étudie une situation contextualisée intitulée \og Réaménagement d'une cour \fg. Les données sont les suivantes : une grandeur de départ vaut $120$, elle évolue selon un coefficient multiplicateur $1{,}08$, puis on lui retranche $15$ unités fixes. Une autre grandeur est donnée par la fonction $f(x)=3x+24$.
\begin{enumerate}
  \item Expliquer pourquoi une augmentation de $8\%$ correspond à une multiplication par $1{,}08$.
  \item Calculer la valeur obtenue après augmentation à partir de $120$.
  \item Retirer ensuite $15$ unités et donner le résultat final.
  \item Calculer $f(0)$, $f(12)$ et $f(25)$.
  \item Déterminer l'antécédent de $99$ par la fonction $f$.
  \item Interpréter ce résultat dans le contexte du sujet.
\end{enumerate}
\end{exercice}

\begin{exercice}[M-LONG-10 -- Partie B : calcul littéral et équation]
On considère l'expression :
\[
E=(2x-5)^2-(x+3)(x-3).
\]
\begin{enumerate}
  \item Développer $(2x-5)^2$.
  \item Développer $(x+3)(x-3)$.
  \item En déduire une forme développée et réduite de $E$.
  \item Calculer $E$ pour $x=4$.
  \item Résoudre l'équation $E=25$.
  \item Vérifier une solution obtenue.
\end{enumerate}
\end{exercice}

\begin{exercice}[M-LONG-10 -- Partie C : géométrie et grandeurs]
Un triangle rectangle est utilisé dans la modélisation. Ses côtés de l'angle droit mesurent $9\ \mathrm{cm}$ et $12\ \mathrm{cm}$.
\begin{enumerate}
  \item Calculer la longueur de l'hypoténuse.
  \item Calculer l'aire du triangle.
  \item Calculer $\cos(\widehat{B})$ pour l'angle adjacent au côté de longueur $9\ \mathrm{cm}$.
  \item En déduire une valeur approchée de l'angle.
  \item On agrandit la figure avec un coefficient $1{,}4$. Calculer les nouvelles longueurs.
  \item Par combien l'aire est-elle multipliée ?
\end{enumerate}
\end{exercice}

\begin{exercice}[M-LONG-10 -- Partie D : synthèse et prise d'initiative]
Une décision doit être prise à partir des résultats des parties précédentes.
\begin{enumerate}
  \item Rédiger une conclusion de cinq lignes qui utilise au moins deux résultats numériques obtenus.
  \item Vérifier la cohérence des unités.
  \item Proposer une amélioration du modèle.
  \item Donner une limite du raisonnement.
  \item Écrire une phrase finale comme dans une copie de brevet.
\end{enumerate}
\end{exercice}
\newpage
\subsection{M-LONG-11 -- Batterie externe et usages}

\begin{objectif}
Durée conseillée : 1 h à 1 h 45.\\
Chapitres : fonctions, proportionnalité, puissances de dix, pourcentages.
\end{objectif}
\begin{exercice}[M-LONG-11 -- Partie A : étude de situation]
On étudie une situation contextualisée intitulée \og Batterie externe et usages \fg. Les données sont les suivantes : une grandeur de départ vaut $120$, elle évolue selon un coefficient multiplicateur $1{,}08$, puis on lui retranche $15$ unités fixes. Une autre grandeur est donnée par la fonction $f(x)=3x+24$.
\begin{enumerate}
  \item Expliquer pourquoi une augmentation de $8\%$ correspond à une multiplication par $1{,}08$.
  \item Calculer la valeur obtenue après augmentation à partir de $120$.
  \item Retirer ensuite $15$ unités et donner le résultat final.
  \item Calculer $f(0)$, $f(12)$ et $f(25)$.
  \item Déterminer l'antécédent de $99$ par la fonction $f$.
  \item Interpréter ce résultat dans le contexte du sujet.
\end{enumerate}
\end{exercice}

\begin{exercice}[M-LONG-11 -- Partie B : calcul littéral et équation]
On considère l'expression :
\[
E=(2x-5)^2-(x+3)(x-3).
\]
\begin{enumerate}
  \item Développer $(2x-5)^2$.
  \item Développer $(x+3)(x-3)$.
  \item En déduire une forme développée et réduite de $E$.
  \item Calculer $E$ pour $x=4$.
  \item Résoudre l'équation $E=25$.
  \item Vérifier une solution obtenue.
\end{enumerate}
\end{exercice}

\begin{exercice}[M-LONG-11 -- Partie C : géométrie et grandeurs]
Un triangle rectangle est utilisé dans la modélisation. Ses côtés de l'angle droit mesurent $9\ \mathrm{cm}$ et $12\ \mathrm{cm}$.
\begin{enumerate}
  \item Calculer la longueur de l'hypoténuse.
  \item Calculer l'aire du triangle.
  \item Calculer $\cos(\widehat{B})$ pour l'angle adjacent au côté de longueur $9\ \mathrm{cm}$.
  \item En déduire une valeur approchée de l'angle.
  \item On agrandit la figure avec un coefficient $1{,}4$. Calculer les nouvelles longueurs.
  \item Par combien l'aire est-elle multipliée ?
\end{enumerate}
\end{exercice}

\begin{exercice}[M-LONG-11 -- Partie D : synthèse et prise d'initiative]
Une décision doit être prise à partir des résultats des parties précédentes.
\begin{enumerate}
  \item Rédiger une conclusion de cinq lignes qui utilise au moins deux résultats numériques obtenus.
  \item Vérifier la cohérence des unités.
  \item Proposer une amélioration du modèle.
  \item Donner une limite du raisonnement.
  \item Écrire une phrase finale comme dans une copie de brevet.
\end{enumerate}
\end{exercice}
\newpage
\subsection{M-LONG-12 -- Mini-sujet d'excellence brevet}

\begin{objectif}
Durée conseillée : 1 h à 1 h 45.\\
Chapitres : calculs, fonctions, géométrie, probabilités, algorithmique.
\end{objectif}
\begin{exercice}[M-LONG-12 -- Partie A : étude de situation]
On étudie une situation contextualisée intitulée \og Mini-sujet d'excellence brevet \fg. Les données sont les suivantes : une grandeur de départ vaut $120$, elle évolue selon un coefficient multiplicateur $1{,}08$, puis on lui retranche $15$ unités fixes. Une autre grandeur est donnée par la fonction $f(x)=3x+24$.
\begin{enumerate}
  \item Expliquer pourquoi une augmentation de $8\%$ correspond à une multiplication par $1{,}08$.
  \item Calculer la valeur obtenue après augmentation à partir de $120$.
  \item Retirer ensuite $15$ unités et donner le résultat final.
  \item Calculer $f(0)$, $f(12)$ et $f(25)$.
  \item Déterminer l'antécédent de $99$ par la fonction $f$.
  \item Interpréter ce résultat dans le contexte du sujet.
\end{enumerate}
\end{exercice}

\begin{exercice}[M-LONG-12 -- Partie B : calcul littéral et équation]
On considère l'expression :
\[
E=(2x-5)^2-(x+3)(x-3).
\]
\begin{enumerate}
  \item Développer $(2x-5)^2$.
  \item Développer $(x+3)(x-3)$.
  \item En déduire une forme développée et réduite de $E$.
  \item Calculer $E$ pour $x=4$.
  \item Résoudre l'équation $E=25$.
  \item Vérifier une solution obtenue.
\end{enumerate}
\end{exercice}

\begin{exercice}[M-LONG-12 -- Partie C : géométrie et grandeurs]
Un triangle rectangle est utilisé dans la modélisation. Ses côtés de l'angle droit mesurent $9\ \mathrm{cm}$ et $12\ \mathrm{cm}$.
\begin{enumerate}
  \item Calculer la longueur de l'hypoténuse.
  \item Calculer l'aire du triangle.
  \item Calculer $\cos(\widehat{B})$ pour l'angle adjacent au côté de longueur $9\ \mathrm{cm}$.
  \item En déduire une valeur approchée de l'angle.
  \item On agrandit la figure avec un coefficient $1{,}4$. Calculer les nouvelles longueurs.
  \item Par combien l'aire est-elle multipliée ?
\end{enumerate}
\end{exercice}

\begin{exercice}[M-LONG-12 -- Partie D : synthèse et prise d'initiative]
Une décision doit être prise à partir des résultats des parties précédentes.
\begin{enumerate}
  \item Rédiger une conclusion de cinq lignes qui utilise au moins deux résultats numériques obtenus.
  \item Vérifier la cohérence des unités.
  \item Proposer une amélioration du modèle.
  \item Donner une limite du raisonnement.
  \item Écrire une phrase finale comme dans une copie de brevet.
\end{enumerate}
\end{exercice}
\newpage
\section{Sujets inventés longs de physique-chimie}

\begin{attention}
Les sujets de physique-chimie sont conçus pour travailler les formules, les conversions, les unités, la rédaction scientifique et l'interprétation des résultats.
\end{attention}
\subsection{PC-LONG-01 -- Vélo électrique, énergie et vitesse}

\begin{objectif}
Durée conseillée : 1 h.\\
Compétences : identifier les données, choisir une formule, convertir les unités, conclure scientifiquement.
\end{objectif}
\begin{exercice}[PC-LONG-01 -- Partie A : vitesse]
Un élève utilise un vélo électrique pour parcourir $18$ km. Il met $54$ minutes.
\begin{enumerate}
  \item Convertir $54$ minutes en heures.
  \item Calculer la vitesse moyenne en $\mathrm{km/h}$.
  \item Convertir cette vitesse en $\mathrm{m/s}$.
  \item La vitesse maximale autorisée avec assistance est $25\ \mathrm{km/h}$. Le vélo respecte-t-il cette limite en moyenne ?
\end{enumerate}
\end{exercice}

\begin{exercice}[PC-LONG-01 -- Partie B : énergie consommée]
La batterie a une capacité de $500\ \mathrm{Wh}$. Pendant le trajet, le moteur fournit une puissance moyenne de $180\ \mathrm{W}$.
\begin{enumerate}
  \item Rappeler la formule reliant énergie, puissance et durée.
  \item Calculer l'énergie consommée pendant le trajet.
  \item Quelle fraction de la batterie est utilisée ?
  \item En déduire si la batterie suffit pour faire l'aller-retour.
\end{enumerate}
\end{exercice}

\begin{exercice}[PC-LONG-01 -- Partie C : coût de recharge]
Le prix de l'électricité est $0{,}25$ euro par kWh.
\begin{enumerate}
  \item Convertir $500\ \mathrm{Wh}$ en $\mathrm{kWh}$.
  \item Calculer le coût d'une recharge complète.
  \item Comparer ce coût à un ticket de bus à $1{,}80$ euro.
  \item Donner une limite de cette comparaison.
\end{enumerate}
\end{exercice}
\newpage
\subsection{PC-LONG-02 -- Qualité de l'eau, masse volumique et ions}

\begin{objectif}
Durée conseillée : 1 h.\\
Compétences : identifier les données, choisir une formule, convertir les unités, conclure scientifiquement.
\end{objectif}
\begin{exercice}[PC-LONG-02 -- Partie A : masse volumique]
On prélève un échantillon d'eau salée. Un volume de $250\ \mathrm{mL}$ a une masse de $260\ \mathrm{g}$.
\begin{enumerate}
  \item Convertir $250\ \mathrm{mL}$ en $\mathrm{cm^3}$.
  \item Calculer la masse volumique en $\mathrm{g/cm^3}$.
  \item Comparer avec l'eau pure, de masse volumique environ $1{,}0\ \mathrm{g/cm^3}$.
  \item Proposer une explication.
\end{enumerate}
\end{exercice}

\begin{exercice}[PC-LONG-02 -- Partie B : ions]
On réalise des tests chimiques sur cette eau. Avec du nitrate d'argent, on observe un précipité blanc. Avec de la soude, on n'observe pas de précipité coloré.
\begin{enumerate}
  \item Quel ion peut être mis en évidence par le nitrate d'argent ?
  \item Que signifie l'apparition d'un précipité ?
  \item Pourquoi le résultat avec la soude ne permet-il pas d'affirmer la présence d'ions cuivre ou fer ?
  \item Donner une conclusion sur la composition probable de l'eau.
\end{enumerate}
\end{exercice}

\begin{exercice}[PC-LONG-02 -- Partie C : évaporation]
On laisse évaporer $100\ \mathrm{mL}$ de cette eau salée. Il reste $3{,}5\ \mathrm{g}$ de solide.
\begin{enumerate}
  \item Quelle masse de solide obtiendrait-on avec $1\ \mathrm{L}$ d'eau salée ?
  \item Exprimer cette concentration massique en $\mathrm{g/L}$.
  \item Peut-on affirmer que le solide est uniquement du sel ? Justifier.
  \item Donner deux sources d'erreur possibles dans l'expérience.
\end{enumerate}
\end{exercice}
\newpage
\subsection{PC-LONG-03 -- Électricité domestique et coût}

\begin{objectif}
Durée conseillée : 1 h.\\
Compétences : identifier les données, choisir une formule, convertir les unités, conclure scientifiquement.
\end{objectif}
\begin{exercice}[PC-LONG-03 -- Partie A : données et conversions]
Dans le problème \og Électricité domestique et coût \fg, on mesure une durée de $36$ minutes, une distance de $12$ km et une puissance de $250$ W.
\begin{enumerate}
  \item Convertir $36$ minutes en heures.
  \item Convertir $12$ km en mètres.
  \item Convertir $250$ W en kW.
  \item Calculer la vitesse moyenne en $\mathrm{km/h}$.
  \item Convertir cette vitesse en $\mathrm{m/s}$.
\end{enumerate}
\end{exercice}

\begin{exercice}[PC-LONG-03 -- Partie B : formule physique]
On utilise la relation $E=P	imes t$.
\begin{enumerate}
  \item Préciser les unités possibles de $E$, $P$ et $t$.
  \item Calculer l'énergie consommée si l'appareil fonctionne $36$ minutes.
  \item Exprimer le résultat en Wh puis en kWh.
  \item Calculer le coût pour un prix de $0{,}25$ euro par kWh.
  \item Rédiger une phrase de conclusion.
\end{enumerate}
\end{exercice}

\begin{exercice}[PC-LONG-03 -- Partie C : analyse critique]
\begin{enumerate}
  \item Citer deux sources d'erreur possibles dans les mesures.
  \item Expliquer pourquoi les unités doivent être vérifiées.
  \item Proposer une amélioration du protocole.
  \item Dire si le résultat obtenu paraît plausible.
  \item Conclure en langage scientifique.
\end{enumerate}
\end{exercice}
\newpage
\subsection{PC-LONG-04 -- Sonar, écho et signaux}

\begin{objectif}
Durée conseillée : 1 h.\\
Compétences : identifier les données, choisir une formule, convertir les unités, conclure scientifiquement.
\end{objectif}
\begin{exercice}[PC-LONG-04 -- Partie A : données et conversions]
Dans le problème \og Sonar, écho et signaux \fg, on mesure une durée de $36$ minutes, une distance de $12$ km et une puissance de $250$ W.
\begin{enumerate}
  \item Convertir $36$ minutes en heures.
  \item Convertir $12$ km en mètres.
  \item Convertir $250$ W en kW.
  \item Calculer la vitesse moyenne en $\mathrm{km/h}$.
  \item Convertir cette vitesse en $\mathrm{m/s}$.
\end{enumerate}
\end{exercice}

\begin{exercice}[PC-LONG-04 -- Partie B : formule physique]
On utilise la relation $E=P	imes t$.
\begin{enumerate}
  \item Préciser les unités possibles de $E$, $P$ et $t$.
  \item Calculer l'énergie consommée si l'appareil fonctionne $36$ minutes.
  \item Exprimer le résultat en Wh puis en kWh.
  \item Calculer le coût pour un prix de $0{,}25$ euro par kWh.
  \item Rédiger une phrase de conclusion.
\end{enumerate}
\end{exercice}

\begin{exercice}[PC-LONG-04 -- Partie C : analyse critique]
\begin{enumerate}
  \item Citer deux sources d'erreur possibles dans les mesures.
  \item Expliquer pourquoi les unités doivent être vérifiées.
  \item Proposer une amélioration du protocole.
  \item Dire si le résultat obtenu paraît plausible.
  \item Conclure en langage scientifique.
\end{enumerate}
\end{exercice}
\newpage
\subsection{PC-LONG-05 -- Transformation chimique et conservation des atomes}

\begin{objectif}
Durée conseillée : 1 h.\\
Compétences : identifier les données, choisir une formule, convertir les unités, conclure scientifiquement.
\end{objectif}
\begin{exercice}[PC-LONG-05 -- Partie A : données et conversions]
Dans le problème \og Transformation chimique et conservation des atomes \fg, on mesure une durée de $36$ minutes, une distance de $12$ km et une puissance de $250$ W.
\begin{enumerate}
  \item Convertir $36$ minutes en heures.
  \item Convertir $12$ km en mètres.
  \item Convertir $250$ W en kW.
  \item Calculer la vitesse moyenne en $\mathrm{km/h}$.
  \item Convertir cette vitesse en $\mathrm{m/s}$.
\end{enumerate}
\end{exercice}

\begin{exercice}[PC-LONG-05 -- Partie B : formule physique]
On utilise la relation $E=P	imes t$.
\begin{enumerate}
  \item Préciser les unités possibles de $E$, $P$ et $t$.
  \item Calculer l'énergie consommée si l'appareil fonctionne $36$ minutes.
  \item Exprimer le résultat en Wh puis en kWh.
  \item Calculer le coût pour un prix de $0{,}25$ euro par kWh.
  \item Rédiger une phrase de conclusion.
\end{enumerate}
\end{exercice}

\begin{exercice}[PC-LONG-05 -- Partie C : analyse critique]
\begin{enumerate}
  \item Citer deux sources d'erreur possibles dans les mesures.
  \item Expliquer pourquoi les unités doivent être vérifiées.
  \item Proposer une amélioration du protocole.
  \item Dire si le résultat obtenu paraît plausible.
  \item Conclure en langage scientifique.
\end{enumerate}
\end{exercice}
\newpage
\subsection{PC-LONG-06 -- Isolation thermique simplifiée}

\begin{objectif}
Durée conseillée : 1 h.\\
Compétences : identifier les données, choisir une formule, convertir les unités, conclure scientifiquement.
\end{objectif}
\begin{exercice}[PC-LONG-06 -- Partie A : données et conversions]
Dans le problème \og Isolation thermique simplifiée \fg, on mesure une durée de $36$ minutes, une distance de $12$ km et une puissance de $250$ W.
\begin{enumerate}
  \item Convertir $36$ minutes en heures.
  \item Convertir $12$ km en mètres.
  \item Convertir $250$ W en kW.
  \item Calculer la vitesse moyenne en $\mathrm{km/h}$.
  \item Convertir cette vitesse en $\mathrm{m/s}$.
\end{enumerate}
\end{exercice}

\begin{exercice}[PC-LONG-06 -- Partie B : formule physique]
On utilise la relation $E=P	imes t$.
\begin{enumerate}
  \item Préciser les unités possibles de $E$, $P$ et $t$.
  \item Calculer l'énergie consommée si l'appareil fonctionne $36$ minutes.
  \item Exprimer le résultat en Wh puis en kWh.
  \item Calculer le coût pour un prix de $0{,}25$ euro par kWh.
  \item Rédiger une phrase de conclusion.
\end{enumerate}
\end{exercice}

\begin{exercice}[PC-LONG-06 -- Partie C : analyse critique]
\begin{enumerate}
  \item Citer deux sources d'erreur possibles dans les mesures.
  \item Expliquer pourquoi les unités doivent être vérifiées.
  \item Proposer une amélioration du protocole.
  \item Dire si le résultat obtenu paraît plausible.
  \item Conclure en langage scientifique.
\end{enumerate}
\end{exercice}
\newpage
\section{Brevet blanc inventé complet}

\begin{objectif}
Durée conseillée : 2 h. Sujet inventé, un peu plus difficile qu'un brevet standard. Il doit être fait sans aide, puis corrigé soigneusement.
\end{objectif}

\begin{exercice}[BB-MATHS-01 -- Exercice 1 : calculs]
Calculer en détaillant :
\[
A=5-2(3-7)^2+\frac{3}{4}\div\frac{9}{8}.
\]
\[
B=\frac{5}{6}+\frac{7}{9}\times\frac{3}{14}.
\]
\[
C=\frac{3^5\times3^{-2}}{3}.
\]
\end{exercice}

\begin{exercice}[BB-MATHS-01 -- Exercice 2 : calcul littéral]
On considère :
\[
E=(2x-3)^2-(x+5)(x-5).
\]
\begin{enumerate}
  \item Développer et réduire $E$.
  \item Calculer $E$ pour $x=2$.
  \item Résoudre l'équation $E=34$.
\end{enumerate}
\end{exercice}

\begin{exercice}[BB-MATHS-01 -- Exercice 3 : fonctions]
Une salle de sport propose deux abonnements : formule A, $18$ euros par séance ; formule B, $80$ euros d'inscription puis $10$ euros par séance. On note $x$ le nombre de séances.
\begin{enumerate}
  \item Écrire $A(x)$ et $B(x)$.
  \item Calculer $A(8)$ et $B(8)$.
  \item Résoudre $B(x)<A(x)$.
  \item Interpréter la réponse.
\end{enumerate}
\end{exercice}

\begin{exercice}[BB-MATHS-01 -- Exercice 4 : géométrie]
Un triangle $ABC$ vérifie :
\[
AB=28\ \mathrm{cm},\quad AC=45\ \mathrm{cm},\quad BC=53\ \mathrm{cm}.
\]
\begin{enumerate}
  \item Montrer que le triangle est rectangle.
  \item Calculer son aire.
  \item Calculer $\sin(\widehat{B})$.
  \item En déduire une valeur approchée de $\widehat{B}$.
\end{enumerate}
\end{exercice}

\begin{exercice}[BB-MATHS-01 -- Exercice 5 : statistiques et probabilités]
Dans un club, les âges des 12 adhérents sont :
\[
11,12,12,13,14,14,14,15,16,16,17,18.
\]
\begin{enumerate}
  \item Calculer l'étendue.
  \item Calculer la moyenne.
  \item Donner la médiane.
  \item On choisit un adhérent au hasard. Quelle est la probabilité qu'il ait au moins $15$ ans ?
\end{enumerate}
\end{exercice}
\newpage
\section{Banque d'automatismes}

\begin{objectif}
Ces exercices courts servent à sécuriser les automatismes. Ils peuvent être faits sans calculatrice pour une partie d'entre eux.
\end{objectif}
\begin{exercice}[Automatismes série 1]
Calculer ou répondre rapidement.
\begin{multicols}{2}
\begin{enumerate}[label=\alph*)]
  \item $7+3	imes 5$
  \item $(7+3)	imes 5$
  \item $2^4+3^2$
  \item $\dfrac{2}{3}+\dfrac{5}{6}$
  \item $\dfrac{4}{7}	imes\dfrac{3}{5}$
  \item $\sqrt{81}+\sqrt{49}$
  \item $3x+5x-2x$
  \item $4(x-3)$
  \item Résoudre $2x+5=17$.
  \item Donner $15\%$ de $80$.
\end{enumerate}
\end{multicols}
\end{exercice}
\begin{exercice}[Automatismes série 2]
Calculer ou répondre rapidement.
\begin{multicols}{2}
\begin{enumerate}[label=\alph*)]
  \item $7+3	imes 5$
  \item $(7+3)	imes 5$
  \item $2^4+3^2$
  \item $\dfrac{2}{3}+\dfrac{5}{6}$
  \item $\dfrac{4}{7}	imes\dfrac{3}{5}$
  \item $\sqrt{81}+\sqrt{49}$
  \item $3x+5x-2x$
  \item $4(x-3)$
  \item Résoudre $2x+5=17$.
  \item Donner $15\%$ de $80$.
\end{enumerate}
\end{multicols}
\end{exercice}
\begin{exercice}[Automatismes série 3]
Calculer ou répondre rapidement.
\begin{multicols}{2}
\begin{enumerate}[label=\alph*)]
  \item $7+3	imes 5$
  \item $(7+3)	imes 5$
  \item $2^4+3^2$
  \item $\dfrac{2}{3}+\dfrac{5}{6}$
  \item $\dfrac{4}{7}	imes\dfrac{3}{5}$
  \item $\sqrt{81}+\sqrt{49}$
  \item $3x+5x-2x$
  \item $4(x-3)$
  \item Résoudre $2x+5=17$.
  \item Donner $15\%$ de $80$.
\end{enumerate}
\end{multicols}
\end{exercice}
\begin{exercice}[Automatismes série 4]
Calculer ou répondre rapidement.
\begin{multicols}{2}
\begin{enumerate}[label=\alph*)]
  \item $7+3	imes 5$
  \item $(7+3)	imes 5$
  \item $2^4+3^2$
  \item $\dfrac{2}{3}+\dfrac{5}{6}$
  \item $\dfrac{4}{7}	imes\dfrac{3}{5}$
  \item $\sqrt{81}+\sqrt{49}$
  \item $3x+5x-2x$
  \item $4(x-3)$
  \item Résoudre $2x+5=17$.
  \item Donner $15\%$ de $80$.
\end{enumerate}
\end{multicols}
\end{exercice}
\begin{exercice}[Automatismes série 5]
Calculer ou répondre rapidement.
\begin{multicols}{2}
\begin{enumerate}[label=\alph*)]
  \item $7+3	imes 5$
  \item $(7+3)	imes 5$
  \item $2^4+3^2$
  \item $\dfrac{2}{3}+\dfrac{5}{6}$
  \item $\dfrac{4}{7}	imes\dfrac{3}{5}$
  \item $\sqrt{81}+\sqrt{49}$
  \item $3x+5x-2x$
  \item $4(x-3)$
  \item Résoudre $2x+5=17$.
  \item Donner $15\%$ de $80$.
\end{enumerate}
\end{multicols}
\end{exercice}
\begin{exercice}[Automatismes série 6]
Calculer ou répondre rapidement.
\begin{multicols}{2}
\begin{enumerate}[label=\alph*)]
  \item $7+3	imes 5$
  \item $(7+3)	imes 5$
  \item $2^4+3^2$
  \item $\dfrac{2}{3}+\dfrac{5}{6}$
  \item $\dfrac{4}{7}	imes\dfrac{3}{5}$
  \item $\sqrt{81}+\sqrt{49}$
  \item $3x+5x-2x$
  \item $4(x-3)$
  \item Résoudre $2x+5=17$.
  \item Donner $15\%$ de $80$.
\end{enumerate}
\end{multicols}
\end{exercice}
\begin{exercice}[Automatismes série 7]
Calculer ou répondre rapidement.
\begin{multicols}{2}
\begin{enumerate}[label=\alph*)]
  \item $7+3	imes 5$
  \item $(7+3)	imes 5$
  \item $2^4+3^2$
  \item $\dfrac{2}{3}+\dfrac{5}{6}$
  \item $\dfrac{4}{7}	imes\dfrac{3}{5}$
  \item $\sqrt{81}+\sqrt{49}$
  \item $3x+5x-2x$
  \item $4(x-3)$
  \item Résoudre $2x+5=17$.
  \item Donner $15\%$ de $80$.
\end{enumerate}
\end{multicols}
\end{exercice}
\begin{exercice}[Automatismes série 8]
Calculer ou répondre rapidement.
\begin{multicols}{2}
\begin{enumerate}[label=\alph*)]
  \item $7+3	imes 5$
  \item $(7+3)	imes 5$
  \item $2^4+3^2$
  \item $\dfrac{2}{3}+\dfrac{5}{6}$
  \item $\dfrac{4}{7}	imes\dfrac{3}{5}$
  \item $\sqrt{81}+\sqrt{49}$
  \item $3x+5x-2x$
  \item $4(x-3)$
  \item Résoudre $2x+5=17$.
  \item Donner $15\%$ de $80$.
\end{enumerate}
\end{multicols}
\end{exercice}
\begin{exercice}[Automatismes série 9]
Calculer ou répondre rapidement.
\begin{multicols}{2}
\begin{enumerate}[label=\alph*)]
  \item $7+3	imes 5$
  \item $(7+3)	imes 5$
  \item $2^4+3^2$
  \item $\dfrac{2}{3}+\dfrac{5}{6}$
  \item $\dfrac{4}{7}	imes\dfrac{3}{5}$
  \item $\sqrt{81}+\sqrt{49}$
  \item $3x+5x-2x$
  \item $4(x-3)$
  \item Résoudre $2x+5=17$.
  \item Donner $15\%$ de $80$.
\end{enumerate}
\end{multicols}
\end{exercice}
\begin{exercice}[Automatismes série 10]
Calculer ou répondre rapidement.
\begin{multicols}{2}
\begin{enumerate}[label=\alph*)]
  \item $7+3	imes 5$
  \item $(7+3)	imes 5$
  \item $2^4+3^2$
  \item $\dfrac{2}{3}+\dfrac{5}{6}$
  \item $\dfrac{4}{7}	imes\dfrac{3}{5}$
  \item $\sqrt{81}+\sqrt{49}$
  \item $3x+5x-2x$
  \item $4(x-3)$
  \item Résoudre $2x+5=17$.
  \item Donner $15\%$ de $80$.
\end{enumerate}
\end{multicols}
\end{exercice}
\begin{exercice}[Automatismes série 11]
Calculer ou répondre rapidement.
\begin{multicols}{2}
\begin{enumerate}[label=\alph*)]
  \item $7+3	imes 5$
  \item $(7+3)	imes 5$
  \item $2^4+3^2$
  \item $\dfrac{2}{3}+\dfrac{5}{6}$
  \item $\dfrac{4}{7}	imes\dfrac{3}{5}$
  \item $\sqrt{81}+\sqrt{49}$
  \item $3x+5x-2x$
  \item $4(x-3)$
  \item Résoudre $2x+5=17$.
  \item Donner $15\%$ de $80$.
\end{enumerate}
\end{multicols}
\end{exercice}
\begin{exercice}[Automatismes série 12]
Calculer ou répondre rapidement.
\begin{multicols}{2}
\begin{enumerate}[label=\alph*)]
  \item $7+3	imes 5$
  \item $(7+3)	imes 5$
  \item $2^4+3^2$
  \item $\dfrac{2}{3}+\dfrac{5}{6}$
  \item $\dfrac{4}{7}	imes\dfrac{3}{5}$
  \item $\sqrt{81}+\sqrt{49}$
  \item $3x+5x-2x$
  \item $4(x-3)$
  \item Résoudre $2x+5=17$.
  \item Donner $15\%$ de $80$.
\end{enumerate}
\end{multicols}
\end{exercice}
\begin{exercice}[Automatismes série 13]
Calculer ou répondre rapidement.
\begin{multicols}{2}
\begin{enumerate}[label=\alph*)]
  \item $7+3	imes 5$
  \item $(7+3)	imes 5$
  \item $2^4+3^2$
  \item $\dfrac{2}{3}+\dfrac{5}{6}$
  \item $\dfrac{4}{7}	imes\dfrac{3}{5}$
  \item $\sqrt{81}+\sqrt{49}$
  \item $3x+5x-2x$
  \item $4(x-3)$
  \item Résoudre $2x+5=17$.
  \item Donner $15\%$ de $80$.
\end{enumerate}
\end{multicols}
\end{exercice}
\begin{exercice}[Automatismes série 14]
Calculer ou répondre rapidement.
\begin{multicols}{2}
\begin{enumerate}[label=\alph*)]
  \item $7+3	imes 5$
  \item $(7+3)	imes 5$
  \item $2^4+3^2$
  \item $\dfrac{2}{3}+\dfrac{5}{6}$
  \item $\dfrac{4}{7}	imes\dfrac{3}{5}$
  \item $\sqrt{81}+\sqrt{49}$
  \item $3x+5x-2x$
  \item $4(x-3)$
  \item Résoudre $2x+5=17$.
  \item Donner $15\%$ de $80$.
\end{enumerate}
\end{multicols}
\end{exercice}
\begin{exercice}[Automatismes série 15]
Calculer ou répondre rapidement.
\begin{multicols}{2}
\begin{enumerate}[label=\alph*)]
  \item $7+3	imes 5$
  \item $(7+3)	imes 5$
  \item $2^4+3^2$
  \item $\dfrac{2}{3}+\dfrac{5}{6}$
  \item $\dfrac{4}{7}	imes\dfrac{3}{5}$
  \item $\sqrt{81}+\sqrt{49}$
  \item $3x+5x-2x$
  \item $4(x-3)$
  \item Résoudre $2x+5=17$.
  \item Donner $15\%$ de $80$.
\end{enumerate}
\end{multicols}
\end{exercice}
\newpage
\section{Corrections détaillées et méthodes de correction active}

\begin{methode}
Pour chaque correction, ne pas se contenter de lire. Il faut refaire le calcul sur une feuille, puis comparer. Si l'erreur revient deux fois, elle doit être inscrite dans la fiche personnelle.
\end{methode}
\begin{correction}[M-LONG-01]
Partie A :
\[
(2-5)^2=(-3)^2=9.
\]
\[
\frac{5}{6}\div\frac{10}{9}=\frac{5}{6}\times\frac{9}{10}=\frac{45}{60}=\frac34.
\]
Donc :
\[
A=7-3\times9+\frac34=7-27+\frac34=-20+\frac34=-\frac{77}{4}.
\]
Pour $B$ :
\[
2^5\times2^{-3}=2^{5-3}=2^2=4.
\]
Donc :
\[
B=\frac44+\frac35-\frac7{10}=1+\frac6{10}-\frac7{10}=\frac9{10}.
\]
Pour $C$ :
\[
\frac34-\frac56=\frac9{12}-\frac{10}{12}=-\frac1{12}.
\]
Donc :
\[
C=\left(-\frac1{12}\right)^2=\frac1{144}.
\]
Partie B : avec $x$ :
\[
x\rightarrow x+4\rightarrow x(x+4)=x^2+4x.
\]
On soustrait $x^2$ puis on ajoute $7$ :
\[
x^2+4x-x^2+7=4x+7.
\]
Pour obtenir $51$ :
\[
4x+7=51\Longleftrightarrow 4x=44\Longleftrightarrow x=11.
\]
Partie C :
\[
E=(x+4)^2-(x+4)(2x-1)=(x+4)[(x+4)-(2x-1)].
\]
\[
E=(x+4)(5-x).
\]
Donc :
\[
E=0\Longleftrightarrow x=-4\ \text{ou}\ x=5.
\]
\end{correction}

\begin{correction}[M-LONG-02]
\[
A(x)=0{,}30x,\qquad B(x)=5+0{,}12x.
\]
Pour $20$ minutes :
\[
A(20)=6,\qquad B(20)=7{,}4.
\]
Pour $45$ minutes :
\[
A(45)=13{,}5,\qquad B(45)=10{,}4.
\]
On résout :
\[
5+0{,}12x<0{,}30x \Longleftrightarrow 5<0{,}18x.
\]
\[
x>\frac{5}{0{,}18}\approx27{,}78.
\]
Le tarif B devient plus avantageux à partir de $28$ minutes. Après augmentation :
\[
A_2(x)=0{,}33x,\qquad B_2(x)=6+0{,}126x.
\]
\end{correction}

\begin{correction}[M-LONG-03]
L'étendue vaut $20-8=12$. La somme des valeurs vaut $211$, donc la moyenne vaut :
\[
\frac{211}{15}\approx14{,}1.
\]
Valeurs rangées :
\[
8,9,10,11,12,13,14,14,15,15,16,17,18,19,20.
\]
La médiane est la 8e valeur, donc $14$. Il y a 7 élèves ayant réussi au moins 15 questions, donc la fréquence est $\frac{7}{15}$.
Pour les probabilités, il y a $14$ questions. Ainsi :
\[
P(\text{géométrie})=\frac{4}{14}=\frac27,
\quad
P(\text{pas physique})=\frac{12}{14}=\frac67.
\]
Deux calculs avec remise :
\[
\frac{5}{14}\times\frac{5}{14}=\frac{25}{196}.
\]
Au moins une question de physique :
\[
1-\left(\frac{12}{14}\right)^2=1-\frac{36}{49}=\frac{13}{49}.
\]
\end{correction}

\begin{correction}[M-LONG-04]
Pour la rampe :
\[
L^2=4{,}80^2+0{,}90^2=23{,}04+0{,}81=23{,}85.
\]
\[
L=\sqrt{23{,}85}\approx4{,}88\ \mathrm{m}.
\]
Pour l'angle $\alpha$ :
\[
\tan(\alpha)=\frac{0{,}90}{4{,}80}=0{,}1875.
\]
Donc $\alpha\approx10{,}6^\circ$. La rampe respecte la norme. Pour le cylindre :
\[
V=\pi r^2h=\pi\times0{,}75^2\times1{,}80\approx3{,}18\ \mathrm{m^3}.
\]
Donc $V\approx3180\ \mathrm{L}$. À $80\%$, le volume vaut $2544\ \mathrm{L}$. Avec $150$ L par jour :
\[
\frac{2544}{150}\approx16{,}96.
\]
La famille peut tenir environ 16 jours complets. Pour la maquette, les volumes sont multipliés par $0{,}2^3=0{,}008$.
\end{correction}

\begin{correction}[M-LONG-05]
\[
C(x)=720+38x+12x=720+50x.
\]
\[
C(40)=720+50\times40=2720.
\]
Le coût par élève est $2720/40=68$ euros. Condition :
\[
\frac{720+50x}{x}\leq70 \Longleftrightarrow 720+50x\leq70x.
\]
\[
720\leq20x \Longleftrightarrow x\geq36.
\]
Trajet : $90\times2=180$ km, il reste $135$ km. La deuxième partie dure $135/75=1{,}8$ h. Durée totale : $3{,}8$ h. Vitesse moyenne :
\[
\frac{315}{3{,}8}\approx82{,}9\ \mathrm{km/h}.
\]
Chambres : $3a+4b=41$. Deux solutions sont $(a,b)=(3,8)$ et $(7,5)$. La première utilise $11$ chambres, la seconde $12$ chambres. Probabilités :
\[
P(\text{fille})=\frac{24}{41},\quad P(\text{déjà visité})=\frac{10}{41},\quad P(\text{fille et déjà visité})=\frac{6}{41}.
\]
Garçons n'ayant jamais visité : $13$, donc probabilité $13/41$.
\end{correction}
\begin{correction}[M-LONG-06 -- Correction type]
On applique les méthodes classiques : coefficient multiplicateur, fonction affine, développement, Pythagore et trigonométrie. Pour la partie A, une augmentation de $8\%$ correspond à $1{,}08$, donc la valeur $120$ devient $129{,}6$, puis $114{,}6$ après retrait de $15$. Pour la fonction $f(x)=3x+24$, on obtient $f(0)=24$, $f(12)=60$ et $f(25)=99$. L'antécédent de $99$ est donc $25$.

Pour la partie B :
\[
(2x-5)^2=4x^2-20x+25,
\qquad
(x+3)(x-3)=x^2-9.
\]
Donc :
\[
E=4x^2-20x+25-(x^2-9)=3x^2-20x+34.
\]
Pour la partie C, l'hypoténuse vaut :
\[
\sqrt{9^2+12^2}=\sqrt{225}=15.
\]
L'aire vaut $54\ \mathrm{cm^2}$. Si l'agrandissement a pour coefficient $1{,}4$, les aires sont multipliées par $1{,}4^2=1{,}96$.
\end{correction}
\begin{correction}[M-LONG-07 -- Correction type]
On applique les méthodes classiques : coefficient multiplicateur, fonction affine, développement, Pythagore et trigonométrie. Pour la partie A, une augmentation de $8\%$ correspond à $1{,}08$, donc la valeur $120$ devient $129{,}6$, puis $114{,}6$ après retrait de $15$. Pour la fonction $f(x)=3x+24$, on obtient $f(0)=24$, $f(12)=60$ et $f(25)=99$. L'antécédent de $99$ est donc $25$.

Pour la partie B :
\[
(2x-5)^2=4x^2-20x+25,
\qquad
(x+3)(x-3)=x^2-9.
\]
Donc :
\[
E=4x^2-20x+25-(x^2-9)=3x^2-20x+34.
\]
Pour la partie C, l'hypoténuse vaut :
\[
\sqrt{9^2+12^2}=\sqrt{225}=15.
\]
L'aire vaut $54\ \mathrm{cm^2}$. Si l'agrandissement a pour coefficient $1{,}4$, les aires sont multipliées par $1{,}4^2=1{,}96$.
\end{correction}
\begin{correction}[M-LONG-08 -- Correction type]
On applique les méthodes classiques : coefficient multiplicateur, fonction affine, développement, Pythagore et trigonométrie. Pour la partie A, une augmentation de $8\%$ correspond à $1{,}08$, donc la valeur $120$ devient $129{,}6$, puis $114{,}6$ après retrait de $15$. Pour la fonction $f(x)=3x+24$, on obtient $f(0)=24$, $f(12)=60$ et $f(25)=99$. L'antécédent de $99$ est donc $25$.

Pour la partie B :
\[
(2x-5)^2=4x^2-20x+25,
\qquad
(x+3)(x-3)=x^2-9.
\]
Donc :
\[
E=4x^2-20x+25-(x^2-9)=3x^2-20x+34.
\]
Pour la partie C, l'hypoténuse vaut :
\[
\sqrt{9^2+12^2}=\sqrt{225}=15.
\]
L'aire vaut $54\ \mathrm{cm^2}$. Si l'agrandissement a pour coefficient $1{,}4$, les aires sont multipliées par $1{,}4^2=1{,}96$.
\end{correction}
\begin{correction}[M-LONG-09 -- Correction type]
On applique les méthodes classiques : coefficient multiplicateur, fonction affine, développement, Pythagore et trigonométrie. Pour la partie A, une augmentation de $8\%$ correspond à $1{,}08$, donc la valeur $120$ devient $129{,}6$, puis $114{,}6$ après retrait de $15$. Pour la fonction $f(x)=3x+24$, on obtient $f(0)=24$, $f(12)=60$ et $f(25)=99$. L'antécédent de $99$ est donc $25$.

Pour la partie B :
\[
(2x-5)^2=4x^2-20x+25,
\qquad
(x+3)(x-3)=x^2-9.
\]
Donc :
\[
E=4x^2-20x+25-(x^2-9)=3x^2-20x+34.
\]
Pour la partie C, l'hypoténuse vaut :
\[
\sqrt{9^2+12^2}=\sqrt{225}=15.
\]
L'aire vaut $54\ \mathrm{cm^2}$. Si l'agrandissement a pour coefficient $1{,}4$, les aires sont multipliées par $1{,}4^2=1{,}96$.
\end{correction}
\begin{correction}[M-LONG-10 -- Correction type]
On applique les méthodes classiques : coefficient multiplicateur, fonction affine, développement, Pythagore et trigonométrie. Pour la partie A, une augmentation de $8\%$ correspond à $1{,}08$, donc la valeur $120$ devient $129{,}6$, puis $114{,}6$ après retrait de $15$. Pour la fonction $f(x)=3x+24$, on obtient $f(0)=24$, $f(12)=60$ et $f(25)=99$. L'antécédent de $99$ est donc $25$.

Pour la partie B :
\[
(2x-5)^2=4x^2-20x+25,
\qquad
(x+3)(x-3)=x^2-9.
\]
Donc :
\[
E=4x^2-20x+25-(x^2-9)=3x^2-20x+34.
\]
Pour la partie C, l'hypoténuse vaut :
\[
\sqrt{9^2+12^2}=\sqrt{225}=15.
\]
L'aire vaut $54\ \mathrm{cm^2}$. Si l'agrandissement a pour coefficient $1{,}4$, les aires sont multipliées par $1{,}4^2=1{,}96$.
\end{correction}
\begin{correction}[M-LONG-11 -- Correction type]
On applique les méthodes classiques : coefficient multiplicateur, fonction affine, développement, Pythagore et trigonométrie. Pour la partie A, une augmentation de $8\%$ correspond à $1{,}08$, donc la valeur $120$ devient $129{,}6$, puis $114{,}6$ après retrait de $15$. Pour la fonction $f(x)=3x+24$, on obtient $f(0)=24$, $f(12)=60$ et $f(25)=99$. L'antécédent de $99$ est donc $25$.

Pour la partie B :
\[
(2x-5)^2=4x^2-20x+25,
\qquad
(x+3)(x-3)=x^2-9.
\]
Donc :
\[
E=4x^2-20x+25-(x^2-9)=3x^2-20x+34.
\]
Pour la partie C, l'hypoténuse vaut :
\[
\sqrt{9^2+12^2}=\sqrt{225}=15.
\]
L'aire vaut $54\ \mathrm{cm^2}$. Si l'agrandissement a pour coefficient $1{,}4$, les aires sont multipliées par $1{,}4^2=1{,}96$.
\end{correction}
\begin{correction}[M-LONG-12 -- Correction type]
On applique les méthodes classiques : coefficient multiplicateur, fonction affine, développement, Pythagore et trigonométrie. Pour la partie A, une augmentation de $8\%$ correspond à $1{,}08$, donc la valeur $120$ devient $129{,}6$, puis $114{,}6$ après retrait de $15$. Pour la fonction $f(x)=3x+24$, on obtient $f(0)=24$, $f(12)=60$ et $f(25)=99$. L'antécédent de $99$ est donc $25$.

Pour la partie B :
\[
(2x-5)^2=4x^2-20x+25,
\qquad
(x+3)(x-3)=x^2-9.
\]
Donc :
\[
E=4x^2-20x+25-(x^2-9)=3x^2-20x+34.
\]
Pour la partie C, l'hypoténuse vaut :
\[
\sqrt{9^2+12^2}=\sqrt{225}=15.
\]
L'aire vaut $54\ \mathrm{cm^2}$. Si l'agrandissement a pour coefficient $1{,}4$, les aires sont multipliées par $1{,}4^2=1{,}96$.
\end{correction}
\begin{correction}[PC-LONG-01]
\[
54\ \mathrm{min}=\frac{54}{60}=0{,}9\ \mathrm{h}.
\]
\[
v=\frac{18}{0{,}9}=20\ \mathrm{km/h}.
\]
\[
20\ \mathrm{km/h}=\frac{20}{3{,}6}\approx5{,}6\ \mathrm{m/s}.
\]
La vitesse moyenne respecte la limite de $25\ \mathrm{km/h}$.

Énergie consommée :
\[
E=P\times t=180\times0{,}9=162\ \mathrm{Wh}.
\]
Fraction de batterie utilisée :
\[
\frac{162}{500}=0{,}324.
\]
Donc environ $32{,}4\%$ de la batterie. Pour un aller-retour identique, on utiliserait $324\ \mathrm{Wh}$ : la batterie suffit.
\[
500\ \mathrm{Wh}=0{,}5\ \mathrm{kWh},
\qquad
0{,}5\times0{,}25=0{,}125\ \text{euro}.
\]
\end{correction}

\begin{correction}[PC-LONG-02]
\[
250\ \mathrm{mL}=250\ \mathrm{cm^3}.
\]
\[
\rho=\frac{m}{V}=\frac{260}{250}=1{,}04\ \mathrm{g/cm^3}.
\]
L'eau salée est plus dense que l'eau pure. Le nitrate d'argent met en évidence les ions chlorure $\mathrm{Cl^-}$ par formation d'un précipité blanc. Avec $100\ \mathrm{mL}$, on obtient $3{,}5$ g de solide. Avec $1\ \mathrm{L}$, on obtiendrait $35$ g. La concentration massique est donc $35\ \mathrm{g/L}$. On ne peut pas affirmer que tout le solide est du sel pur.
\end{correction}
\begin{correction}[PC-LONG-03 -- Correction type]
On convertit d'abord les unités : $36$ minutes correspondent à $0{,}6$ h et à $2160$ s. La vitesse moyenne pour $12$ km en $0{,}6$ h vaut :
\[
v=\frac{12}{0{,}6}=20\ \mathrm{km/h}.
\]
En mètres par seconde :
\[
v=\frac{20}{3{,}6}\approx5{,}6\ \mathrm{m/s}.
\]
La puissance $250$ W vaut $0{,}250$ kW. L'énergie consommée pendant $0{,}6$ h est :
\[
E=0{,}250	imes0{,}6=0{,}150\ \mathrm{kWh}.
\]
Le coût est :
\[
0{,}150	imes0{,}25=0{,}0375\ 	ext{euro}.
\]
Il faut conclure avec une unité et rappeler les limites des mesures.
\end{correction}
\begin{correction}[PC-LONG-04 -- Correction type]
On convertit d'abord les unités : $36$ minutes correspondent à $0{,}6$ h et à $2160$ s. La vitesse moyenne pour $12$ km en $0{,}6$ h vaut :
\[
v=\frac{12}{0{,}6}=20\ \mathrm{km/h}.
\]
En mètres par seconde :
\[
v=\frac{20}{3{,}6}\approx5{,}6\ \mathrm{m/s}.
\]
La puissance $250$ W vaut $0{,}250$ kW. L'énergie consommée pendant $0{,}6$ h est :
\[
E=0{,}250	imes0{,}6=0{,}150\ \mathrm{kWh}.
\]
Le coût est :
\[
0{,}150	imes0{,}25=0{,}0375\ 	ext{euro}.
\]
Il faut conclure avec une unité et rappeler les limites des mesures.
\end{correction}
\begin{correction}[PC-LONG-05 -- Correction type]
On convertit d'abord les unités : $36$ minutes correspondent à $0{,}6$ h et à $2160$ s. La vitesse moyenne pour $12$ km en $0{,}6$ h vaut :
\[
v=\frac{12}{0{,}6}=20\ \mathrm{km/h}.
\]
En mètres par seconde :
\[
v=\frac{20}{3{,}6}\approx5{,}6\ \mathrm{m/s}.
\]
La puissance $250$ W vaut $0{,}250$ kW. L'énergie consommée pendant $0{,}6$ h est :
\[
E=0{,}250	imes0{,}6=0{,}150\ \mathrm{kWh}.
\]
Le coût est :
\[
0{,}150	imes0{,}25=0{,}0375\ 	ext{euro}.
\]
Il faut conclure avec une unité et rappeler les limites des mesures.
\end{correction}
\begin{correction}[PC-LONG-06 -- Correction type]
On convertit d'abord les unités : $36$ minutes correspondent à $0{,}6$ h et à $2160$ s. La vitesse moyenne pour $12$ km en $0{,}6$ h vaut :
\[
v=\frac{12}{0{,}6}=20\ \mathrm{km/h}.
\]
En mètres par seconde :
\[
v=\frac{20}{3{,}6}\approx5{,}6\ \mathrm{m/s}.
\]
La puissance $250$ W vaut $0{,}250$ kW. L'énergie consommée pendant $0{,}6$ h est :
\[
E=0{,}250	imes0{,}6=0{,}150\ \mathrm{kWh}.
\]
Le coût est :
\[
0{,}150	imes0{,}25=0{,}0375\ 	ext{euro}.
\]
Il faut conclure avec une unité et rappeler les limites des mesures.
\end{correction}
\newpage
\section{Checklist finale}

\begin{tcolorbox}[colframe=bleu,colback=blue!3!white,title=La veille de l'épreuve de mathématiques]
\begin{itemize}
  \item Je relis les formules essentielles.
  \item Je refais seulement des exercices déjà corrigés.
  \item Je ne commence pas un nouveau chapitre.
  \item Je prépare ma calculatrice, règle, compas, rapporteur.
  \item Je dors correctement.
\end{itemize}
\end{tcolorbox}

\begin{tcolorbox}[colframe=rouge,colback=red!3!white,title=Pendant l'épreuve]
\begin{itemize}
  \item Lire tout le sujet.
  \item Commencer par les exercices les plus sûrs.
  \item Ne pas bloquer trop longtemps.
  \item Écrire les formules utilisées.
  \item Mettre les unités.
  \item Faire des phrases de conclusion.
  \item Garder du temps pour relire.
\end{itemize}
\end{tcolorbox}

\newpage
\section{Banque supplémentaire de problèmes courts}
\begin{objectif}Cette dernière banque sert à augmenter le volume d’entraînement. Les exercices sont courts mais couvrent tout le programme.\end{objectif}
\begin{exercice}[Problème court 1]
On donne une situation de révision rapide.
\begin{enumerate}
  \item Calculer $\dfrac{3}{6}+\dfrac{5}{12}$.
  \item Développer et réduire $(3x-3)(x+4)$.
  \item Résoudre $3x+5=20$.
  \item Un objet parcourt $11$ km en $21$ minutes. Calculer sa vitesse moyenne en km/h.
  \item Rédiger une phrase de conclusion.
\end{enumerate}
\end{exercice}
\begin{exercice}[Problème court 2]
On donne une situation de révision rapide.
\begin{enumerate}
  \item Calculer $\dfrac{4}{6}+\dfrac{5}{12}$.
  \item Développer et réduire $(4x-3)(x+4)$.
  \item Résoudre $4x+5=23$.
  \item Un objet parcourt $12$ km en $22$ minutes. Calculer sa vitesse moyenne en km/h.
  \item Rédiger une phrase de conclusion.
\end{enumerate}
\end{exercice}
\end{document}